% =====================================================================
%  CARNET D'ENTRAÎNEMENT — CHIMIE — FICHE 3 — THÈME 1
%  Particules subatomiques — NIVEAU DIFFICILE — CORRIGÉ
% =====================================================================
\documentclass[11pt,a4paper]{article}

\usepackage[utf8]{inputenc}
\usepackage[T1]{fontenc}
\usepackage{lmodern}
\usepackage[french]{babel}

\usepackage{amsmath,amssymb,mathtools}
\usepackage{xcolor}
\usepackage{colortbl}
\usepackage{tikz}
\usetikzlibrary{arrows.meta,positioning,calc}
\usepackage[most]{tcolorbox}
\usepackage{enumitem}
\usepackage{booktabs}
\usepackage{array}
\usepackage[a4paper,margin=2.1cm,headheight=26pt,headsep=9pt]{geometry}
\usepackage{fancyhdr}
\usepackage[hidelinks]{hyperref}

\definecolor{primary}{RGB}{124,78,140}
\definecolor{primarydark}{RGB}{74,44,92}
\definecolor{defcol}{RGB}{0,114,178}
\definecolor{thmcol}{RGB}{0,140,120}
\definecolor{excol}{RGB}{0,135,90}
\definecolor{attncol}{RGB}{200,40,55}

\newcommand{\nuc}[3]{\prescript{#1}{#2}{\mathrm{#3}}}
\newcommand{\pp}{p^{+}}
\newcommand{\ee}{e^{-}}
\newcommand{\nz}{n^{0}}
\newcommand{\rep}{\textbf{\color{excol}Réponse : }}

\pagestyle{fancy}
\fancyhf{}
\fancyhead[L]{\small\textcolor{primarydark}{\textbf{Fiche 3 — Thème 1}} : Particules subatomiques}
\fancyhead[R]{\small\textcolor{primarydark}{\textsc{Difficile — Corrigé}}}
\fancyfoot[C]{\small\thepage}
\renewcommand{\headrulewidth}{0.8pt}
\renewcommand{\headrule}{\hbox to\headwidth{\color{primary}\leaders\hrule height \headrulewidth\hfill}}

\tcbset{
  boxbase/.style={enhanced,breakable,boxrule=0.8pt,arc=2pt,
     left=2.6mm,right=2.6mm,top=1.8mm,bottom=1.8mm,
     fonttitle=\bfseries,coltitle=white,toptitle=0.4mm,bottomtitle=0.4mm},
}
\newtcolorbox{bilan}[1][]{boxbase,colback=thmcol!7,colframe=thmcol,
  title={\textsc{À retenir}\ifx\relax#1\relax\relax\else\textmd{ : \itshape #1}\fi}}

\newcounter{exo}
\newcommand{\cor}[1][]{\stepcounter{exo}\par\medskip%
  \noindent{\color{primary}\rule[-2pt]{3.5pt}{15pt}}\ \textbf{\color{primarydark}\large Exercice \theexo}%
  \ \textmd{\normalsize\color{black!55}— corrigé}%
  \ifx\relax#1\relax\relax\else\ \textmd{\normalsize\color{black!65}\itshape (#1)}\fi\par\nopagebreak\smallskip}

\setlist{leftmargin=1.7em,topsep=2pt,itemsep=3pt}

\begin{document}

\begin{tikzpicture}
\node[rectangle,rounded corners=6pt,fill=primarydark,text=white,
      minimum width=\textwidth,minimum height=1.35cm,align=center,inner sep=6pt]
  {\Large\bfseries Thème 1 — Particules subatomiques\\[2pt]
   \normalsize\mdseries Niveau Difficile — Corrigé détaillé};
\end{tikzpicture}

% ---------------------------------------------------------------
\cor[vitesse d'un électron à partir de son énergie cinétique]
\begin{enumerate}[label=\alph*)]
  \item $E_c = 18{,}22\times1{,}66\times10^{-19} \approx 3{,}02\times10^{-18}$ J.
  \item De $E_c = \tfrac12 m v^2$ on tire $v = \sqrt{\dfrac{2E_c}{m}}
        = \sqrt{\dfrac{2\times18{,}22\times1{,}66\times10^{-19}}{9{,}11\times10^{-31}}}$.
  \item On simplifie : $\dfrac{2\times18{,}22}{9{,}11} = \dfrac{36{,}44}{9{,}11} = 4$ \emph{exactement},
        et $\dfrac{10^{-19}}{10^{-31}} = 10^{12}$. Donc
        \[ v = \sqrt{4\times1{,}66\times10^{12}} = 2\sqrt{1{,}66}\times10^{6}
             \approx 2{,}58\times10^{6}\ \text{m/s}. \]
        \rep $\boxed{v \approx 2{,}58\times10^{6}\ \text{m/s}}$.
  \item $\dfrac{v}{c} = \dfrac{2{,}58\times10^{6}}{3{,}0\times10^{8}} \approx 0{,}0086 = 0{,}86\,\%$.
        \rep environ $\mathbf{0{,}86\,\%}$ de la vitesse de la lumière — déjà très rapide !
\end{enumerate}

% ---------------------------------------------------------------
\cor[le chemin inverse : de la vitesse à l'énergie]
\begin{enumerate}[label=\alph*)]
  \item $E_c = \tfrac12 m v^2 = \tfrac12\times9{,}11\times10^{-31}\times(3{,}0\times10^{6})^2$.
        Or $(3{,}0\times10^{6})^2 = 9{,}0\times10^{12}$, donc
        \[ E_c = \tfrac12\times9{,}11\times10^{-31}\times9{,}0\times10^{12}
             = \tfrac12\times8{,}20\times10^{-18} \approx 4{,}10\times10^{-18}\ \text{J}. \]
  \item En eV : $E_c = \dfrac{4{,}10\times10^{-18}}{1{,}66\times10^{-19}}
        = \dfrac{4{,}10}{1{,}66}\times10^{1} \approx 24{,}7\ \text{eV}$.
        \rep $\boxed{E_c \approx 4{,}10\times10^{-18}\ \text{J} \approx 24{,}7\ \text{eV}}$.
\end{enumerate}

% ---------------------------------------------------------------
\cor[une microbille chargée]
\begin{enumerate}[label=\alph*)]
  \item La charge est \textbf{négative}, donc la microbille a \textbf{plus d'électrons que de protons} :
        c'est un \textbf{excès d'électrons}.
  \item $N = \dfrac{1{,}6\times10^{-15}}{1{,}6\times10^{-19}}
        = 1{,}0\times10^{\,-15-(-19)} = 1{,}0\times10^{4}$ charges élémentaires.
  \item \rep $\boxed{10^{4} = 10\,000 \text{ électrons en excès}}$.
\end{enumerate}

% ---------------------------------------------------------------
\cor[électron contre proton : même vitesse, quelle énergie ?]
\begin{enumerate}[label=\alph*)]
  \item À vitesse égale, $v$ se simplifie :
        \[ \frac{E_c(\pp)}{E_c(\ee)} = \frac{\tfrac12 m(\pp)\,v^2}{\tfrac12 m(\ee)\,v^2}
             = \frac{m(\pp)}{m(\ee)}. \]
  \item $\dfrac{m(\pp)}{m(\ee)} = \dfrac{1{,}7\times10^{-27}}{9{,}11\times10^{-31}}
        = \dfrac{1{,}7}{9{,}11}\times10^{4} \approx 1{,}9\times10^{3} \approx 2000$.
        \; \rep c'est le \textbf{proton} qui possède la plus grande énergie cinétique, environ
        $\mathbf{2000}$ fois celle de l'électron (car $E_c \propto m$ à vitesse fixée).
\end{enumerate}

\begin{bilan}
$E_c=\tfrac12 m v^2 \Rightarrow v=\sqrt{2E_c/m}$. \textbf{Toujours convertir eV~$\to$~J}
($1$ eV $=1{,}66\times10^{-19}$ J) avant de mélanger avec des kg et des m/s. Le nombre de charges
élémentaires d'un objet vaut $N = |Q|/e$. À vitesse égale, $E_c \propto m$.
\end{bilan}

\end{document}
